\section{Related work}
\label{related}

\subsection{The dataset and classical approaches}

Vidros et al.~\cite{vidros2017} defined online recruitment fraud as a research
problem, released EMSCAD and classified it with bag-of-words features and standard classifiers.
They worked on a balanced subset, so their results do not show how the models behave at the
natural 4.84\% fraud rate. Later studies mostly stay in the same family: text features plus
hand-built meta-features such as a missing company profile or a missing logo, fed to linear,
tree or ensemble classifiers. Taneja et al.~\cite{taneja2025} tabulate ten EMSCAD
studies from 2017 to 2024 whose best models are random forests, rule learners, extra trees or SVMs,
with 91\% to 99.9\% accuracy. The top of that range needs care: Amaar et
al.~\cite{amaar2022} report 99.9\% accuracy on a resampled test set, which does not
reflect the real class balance.

\subsection{Transformer approaches}

Pretrained transformers such as BERT~\cite{devlin2019} have since been applied to
the same data. Akram et al.~\cite{akram2024} fine-tune BERT and RoBERTa with ten
SMOTE variants and report about 90\% balanced accuracy and recall. Fraud-BERT~\cite{taneja2025}
fine-tunes BERT-base on twelve concatenated text columns and reports F1 0.93 on one 80:20 split,
against 0.89 for an SVM and 0.78 for logistic regression on the same split. The paper reports no
seeds, no validation split and no deduplication, although its authors note that EMSCAD contains
duplicate entries, so it is unclear how much of the gap would survive a stricter protocol. A
cheaper alternative to fine-tuning is to reuse a frozen sentence encoder such as
Sentence-BERT~\cite{reimers2019} and train only a small classifier on its embeddings.

\subsection{Evaluation on imbalanced data}

Two lines of work shape how results on such data should be measured. Oversampling methods such
as SMOTE~\cite{chawla2002} rebalance the training data with synthetic minority
examples, which is legitimate during training but makes test numbers meaningless when applied to
the test set. Saito and Rehmsmeier~\cite{saito2015} show that ROC curves look
optimistic when negatives dominate and that precision-recall curves are more informative, because
they ignore the large number of easy true negatives.

\subsection{Explaining neural text classifiers}

A linear model can be read from its coefficients, but a transformer cannot. Integrated
Gradients~\cite{sundararajan2017} explains a single prediction of a neural network by following
the model's gradients from a neutral reference input to the actual input. It has a useful
guarantee, completeness: the attributions of an input add up exactly to the change in the model's
output between the reference and that input.

\subsection{Research gap}

Taken together, reported EMSCAD results are hard to compare because their protocols differ,
single-split results carry no estimate of their spread, and the gain of a fine-tuned transformer
over a tuned word-count baseline is rarely isolated. This project therefore fixes one protocol,
puts every model through it on the same five splits, and asks what each model relies on.
